\chapter{Implementation contract}
\label{app:implementation-contract}

This annex is normative for the first build. The reader-facing chapters explain
why the product exists and what a person should experience. This annex states
the boundaries an implementation may not invent away. The machine-readable
companion is \path{schemas/implementation_contract.json}; the record schemas
are in \path{schemas/}. If this annex and an implementation disagree, the
implementation stops until the contract is deliberately changed and the locked
fixtures are replayed.

The canonical R1--R29 acceptance register is
Table~\ref{tab:expanded-requirements} in Appendix~\ref{app:requirements}.
This annex supplies the execution contract for that register; it does not
create a second list of requirements.

\section{Scope, precedence, and staged gates}
\label{sec:contract-scope}

The production operation is finished-manuscript humanization: an immutable source
plus an explicit authorization produces a separate child revision. The protected
source core is a prerequisite layer within that operation, not an alternative
product. Greenfield authoring from a brief is legacy and experimental, reachable
only under \texttt{hv legacy}, and its records cannot satisfy a release gate.

\begin{table}[H]
\centering
\small
\begin{tabularx}{\textwidth}{@{}p{0.24\textwidth}L p{0.24\textwidth}@{}}
\toprule
Stage & Required output & Gate \\
\midrule
Protected source core & Humanization brief, immutable source snapshot, canonical
LaTeX build, source map, protected-object comparison, versioned records, and safe
command results. & A second operator replays the core from the frozen source and
records. \\
Frozen concept baseline & Complete span partition, concept inventory, explanation
obligations, dependency order, scaffolding dispositions, independent reverse
coverage review, and a named freeze decision. & No reader-facing span is
uncovered and no classification is unresolved. \\
Source-grounded rewriting & Per-passage replacements with located concept
correspondence and obligation evidence; nothing published on truncation. &
Mutations that drop a concept, mention without explaining, add unsupported
content, or reorder a prerequisite block acceptance. \\
Verification and release & Semantic preflight on the candidate, causal repair,
patch assembly, clean revised and blackline builds, and a fail-closed packet. &
Every never-except invariant passes; there is no generic exception. \\
Product evidence & Benchmark manuscript, held-out manuscript, independent replay,
and named-reader results. & Only this gate authorizes external release. \\
\bottomrule
\end{tabularx}
\caption{Staged gates for the humanization path. Each gate reports one evidence
state; a later state is never inferred from an earlier one.}
\label{tab:contract-scope}
\end{table}

When controls conflict, apply them in this order: correctness, source fidelity,
domain meaning, reader comprehension, then template regularity. A safety or
privacy control overrides all five. Project and format exceptions belong in a
configuration record, never in an unexplained manuscript instruction.

\section{Threat model and trust boundary}
\label{sec:contract-threat-model}

Humanvoice handles data that may be both executable and confidential. The
following inputs are untrusted: the authoring brief and evidence files; the
TeX source tree, bibliography, images, custom packages, and auxiliary files;
fixture documents and corpus imports; model prompts and model output; and
reader-supplied answers or uploaded artifacts. Text that looks like an
instruction inside any of those inputs is data. It has no authority to invoke a
tool, disclose a secret, change a policy, or widen a claim.

The default trust boundary is a local isolated workspace. A run receives a
read-only source snapshot and writes to a separate scratch directory. The
canonical source is never compiled in place, and no output from a parser or
model is executed. Remote access is denied unless a separate policy record
names the fields, provider, purpose, retention, and approving owner.

\begin{table}[H]
\centering
\small
\begin{tabularx}{\textwidth}{@{}p{0.09\textwidth}p{0.40\textwidth}L@{}}
\toprule
Control & Rule & Observable test \\
\midrule
T1 & Compile from a read-only snapshot into an allow-listed scratch directory. &
The source tree is byte-identical after every build and no path outside the
allowlist is written. \\
T2 & Invoke TeX with \texttt{-no-shell-escape}; reject
\texttt{\textbackslash write18},
source executables, and model-generated commands. & A shell-command fixture
cannot create a file or alter the run result. \\
T3 & Deny compiler and parser network access. A remote model or metadata lookup
crosses an explicit approval point. & A network-deny fixture blocks an
unapproved call and records exit code 5. \\
T4 & Send critics only delimited, schema-validated JSON. Treat source text and
model output as data, not instructions. & An injection fixture yields a
located finding or abstention and cannot acquire a tool or secret. \\
T5 & Enforce time, memory, file-size, process-count, and output-size limits. &
An exhausted limit aborts the run with its limit and source hash; it never
falls through to a pass. \\
\bottomrule
\end{tabularx}
\caption{The threat model turns local-first and prompt-injection language into
checks a build team can run.}
\label{tab:contract-threat-controls}
\end{table}

The build wrapper must use \texttt{-halt-on-error} and \texttt{-file-line-error}
in addition to \texttt{-no-shell-escape}. A successful PDF is not evidence that
the source was safe: the trust-boundary checks and the protected comparison are
separate release conditions. A security or policy owner may veto a run and may
not be overruled by a model, an editor, or an engineer.

\section{Runtime and model manifest}
\label{sec:contract-runtime}

The deterministic path is the reference implementation. It consists of the
parser adapter, canonical TeX build, protected comparison, citation-identity
checks, and source-map checks. The proposed local inference runtime is
\texttt{llama.cpp}; model quality is a separate question from runtime
availability.

The following are planning candidates, not an unearned adoption claim:

\begin{table}[H]
\centering
\small
\begin{tabularx}{\textwidth}{@{}p{0.22\textwidth}p{0.25\textwidth}L@{}}
\toprule
Role & Candidate & Condition before it may enter a release gate \\
\midrule
Runtime & \texttt{llama.cpp}, exact commit recorded & Rebuild the runtime and
record compiler, hardware, and license information. \\
Reference model & Qwen3--8B--Instruct, GGUF Q4\_K\_M, exact file hash & Run
the frozen calibration and held-out sets with the exact artifact, tokenizer,
prompt template, and sampling settings. \\
Fallback model & Llama 3.1--8B--Instruct, GGUF Q4\_K\_M, exact file hash &
Use only as a separately calibrated alternative; never silently substitute it
for the reference model. \\
No model & Deterministic protected core & May continue when inference is
unavailable, but model-dependent writing and reconstruction gates remain
abstentions and cannot produce a full-path release. \\
\bottomrule
\end{tabularx}
\caption{Model candidates are pinned experimental inputs, not a claim that a
particular model is already suitable.}
\label{tab:contract-models}
\end{table}

Every inference run records the runtime revision, model file hash, tokenizer
hash, prompt-template hash, sampling parameters, seed, hardware, network
state, latency, memory, and output hash. A change to any of those invalidates
the affected calibration and held-out results. The fixtures must be replayed
before the changed component can return a release-ready finding. A model may
suggest a repair; it may not accept the document.

\section{Operating envelope}
\label{sec:contract-envelope}

The numbers below are operating controls, not measured performance and never
prose targets. They make the parser, scratch space, and inference budget
designable. Exhausting a control pauses or blocks work and requires explicit
continuation authorization; it never authorizes dropping a concept, shortening an
explanation, or releasing an incomplete revision. Measured runs either revise the
controls openly or narrow the supported document class.

\begin{table}[H]
\centering
\small
\begin{tabularx}{\textwidth}{@{}p{0.35\textwidth}p{0.19\textwidth}L@{}}
\toprule
Quantity & Control & Interpretation \\
\midrule
Source tree & 250 MiB & Includes bibliography, images, packages, and auxiliary
files in the snapshot. \\
Rendered document & 1,000 pages & A larger document may use the protected core
after an explicit scope decision; it is never silently truncated. \\
Protected objects & 20,000 & Equations, numbers, citations, tables, quotations,
code blocks, and declared claims. \\
Deterministic preflight & 30 minutes & On the reference machine: eight CPU
cores and 32 GiB RAM. \\
Model-assisted verification & 90 minutes & Includes critic calls and
reconstruction; a timeout is an abstention, never a pass. \\
Repair & 3 cycles per unit & Repeated signatures or alternating parent hashes
stop the loop and name an unresolved author choice. \\
Per-call output estimate & 2,000 words & A pre-call segmentation trigger only. A
unit expected to exceed it is split on concept or dependency boundaries and
resumed; content is never trimmed to fit. \\
Inference & Recorded per run & Token, time, and cost totals enter the cost
account and may pause further calls. \\
\bottomrule
\end{tabularx}
\caption{Operating controls. They are scope controls and cost inputs, not
promises of speed and not limits on the length of an adequate explanation.}
\label{tab:contract-envelope}
\end{table}

\subsection{Historical cost observations}

The figures below come from v1 greenfield-drafting runs in August 2026. They are
retained as cost-sizing evidence only. They were produced under a superseded
contract whose per-unit output ceiling caused semantic loss, and they are not
v2 measurements:

\begin{table}[H]
\centering
\small
\begin{tabularx}{\textwidth}{@{}p{0.30\textwidth}p{0.22\textwidth}L@{}}
\toprule
Quantity & Observed under v1 & v2 treatment \\
\midrule
Document size & 15,642 words & No length control; the source defines the work \\
Token cost per run & \$2.24--\$2.53 & Recorded; may pause further calls \\
Wall-clock time & 639 seconds & Recorded against the time control \\
Input tokens per call & 1,020--16,703 & Sized per unit; excess forces a split \\
Output tokens per call & 663--2,147 & Sized per unit; a truncated call publishes
nothing and is resumed \\
Peak memory & Not measured & Recorded against the memory control \\
\bottomrule
\end{tabularx}
\caption{Historical v1 cost observations (2026-08-29), retained for sizing only.
Cost varied roughly 16$\times$ across units because of source size and context
accumulation. Under v2 these are diagnostics: no figure in this table is a target
or a ceiling on an adequate explanation.}
\label{tab:contract-measured-budgets}
\end{table}

Each call receives one source passage with enough neighboring context to keep
transitions coherent, and returns a replacement plus its concept correspondence.
A whole-document model context is outside the release path. A call that cannot
discharge its obligations is split on concept boundaries and resumed; truncation
publishes no replacement and never permits an autonomous rewrite of the
manuscript.

\section{Record and schema policy}
\label{sec:contract-schemas}

Every record carries \texttt{record\_type}, \texttt{schema\_version},
\texttt{record\_id}, \texttt{run\_id}, and \texttt{created\_at}. The record family
is enumerated in \path{schemas/record_catalog.json} and loaded through one shared
registry, \path{src/humanvoice/schemas.py}, so the runtime cannot hold a weaker
duplicate schema. The semantic records are
\path{schemas/humanization-brief.schema.json},
\path{schemas/source-span.schema.json},
\path{schemas/source-concept.schema.json},
\path{schemas/concept-baseline.schema.json},
\path{schemas/concept-dependency.schema.json},
\path{schemas/explanation-obligation.schema.json},
\path{schemas/scaffolding-disposition.schema.json},
\path{schemas/rewrite-unit.schema.json},
\path{schemas/concept-disposition.schema.json},
\path{schemas/concept-correspondence.schema.json}, and
\path{schemas/semantic-preflight-result.schema.json}. The release path adds
\path{schemas/source-snapshot.schema.json},
\path{schemas/protected-manifest.schema.json},
\path{schemas/finding.schema.json}, \path{schemas/revision.schema.json},
\path{schemas/release-decision.schema.json},
\path{schemas/reader-decision.schema.json},
\path{schemas/runtime-manifest.schema.json}, and
\path{schemas/cli-result.schema.json}. The rights record is
\path{schemas/corpus-item.schema.json}. The v1 authoring brief, argument
blueprint, and preflight-run schemas remain readable as \texttt{legacy\_v1}
audit evidence and cannot satisfy a v2 gate.

The compatibility rule is simple. A minor schema version may add optional
fields; readers preserve unknown extension fields. A major version requires a
migration, a changed contract identifier, and replay of the locked fixtures.
The same source, schema, toolchain, configuration, and runtime manifest must
reproduce deterministic finding identifiers. Stochastic outputs retain their
seed and hashes but are never described as bitwise deterministic.

Unknown fields at the top level are rejected. Additions that do not change the
meaning of an existing record belong under an explicit \texttt{extensions}
object and are preserved by readers. The runtime manifest records the compiler,
runtime revision, model and tokenizer hashes when applicable, prompt-template
hash, sampling and seed, hardware, network policy, resource limits, latency,
peak memory, output hash, and an explicit reason for fields that do not apply.
CLI results always include the ten stable output fields named below, and corpus
rights records include \texttt{consent\_date}, including a null value when no
consent date is applicable.

The humanization brief also requires the intended reader, genre, known
vocabulary, named exemplars, evidence boundary, privacy class, remote-inference
authorization, voice constraints, protected-object policy, the authorizing
action, and explicit unknowns. A frozen baseline requires component hashes and
counts, complete partition evidence, both coverage reviews, empty unresolved
sets, and a named freeze decision. A finding carries a consequence, editorial
question, rule version, critic identity, and author disposition as well as its
location and observation.

\section{Command-line contract}
\label{sec:contract-cli}

The production entry point \texttt{hv humanize} and its resumable phases exchange
versioned JSON records. Each invocation writes one machine-readable result to
standard output and human progress or diagnostics to standard error. The stable
output fields are \texttt{contract\_id}, \texttt{contract\_version},
\texttt{command}, \texttt{status}, \texttt{exit\_code}, \texttt{run\_id},
\texttt{record\_paths}, \texttt{blocking\_reasons}, \texttt{source\_hash},
and \texttt{result\_hash}. Messages, display paths, timings, and suggestions
are advisory and may change without a contract-version bump.

\begin{table}[H]
\centering
\small
\begin{tabularx}{\textwidth}{@{}p{0.19\textwidth}p{0.29\textwidth}L@{}}
\toprule
Exit & Meaning & Required behavior \\
\midrule
0 & Pass or released & The command completed its declared job. \\
1 & Deterministic gate failure & Preserve the source and emit the blocking
finding or build result. \\
2 & Abstention or unresolved policy result & Do not treat absence as pass;
identify the affected gate. \\
3 & Invalid brief, input, schema, or usage & Do not invoke the writer. \\
4 & Internal implementation error & Preserve the parent run and provide an
incident identifier. \\
5 & Security or trust-boundary violation & Stop the run, quarantine its output,
and notify the security/policy owner. \\
\bottomrule
\end{tabularx}
\caption{Stable command outcomes. Exit codes carry coarse control flow; the
JSON result carries the reasons.}
\label{tab:contract-exit-codes}
\end{table}

\section{Repair, exceptions, and release authority}
\label{sec:contract-authority}

\texttt{hv repair} never edits the canonical source in place. It requires a
clean version-control tree or an immutable source snapshot, writes a child
revision under \path{.humanvoice/runs/<run\_id>/revisions/<revision\_id>},
validates that child in a temporary directory, and publishes its manifest only
after the comparison passes. The parent and its decisions are never deleted.
Rollback selects the parent revision. A cycle limit, repeated finding
signature, or alternating parent hashes stops the loop and returns an author
choice; it cannot produce a reader packet.

Ordinary release requires every required gate to pass and the document owner to
accept the protected comparison. An exception may be granted by the sponsor or
named project owner. A meaning exception also requires the domain editor's
concurrence; a processing exception requires the security/policy owner's
concurrence. The security/policy owner can veto release.

No exception may waive unresolved protected-object correspondence, an unparsed
object promised by the brief, unauthorized external transmission, missing
critical evidence for a load-bearing claim, or a broken or unreproducible
source build. Every permitted exception records its scope, reason, residual
risk, owner, concurrences, expiry, reader disclosure, and count in the
feasibility record.

\section{Corpus rights and definition of done}
\label{sec:contract-rights-done}

\subsection{Role assignments and coverage}

The R1--R24 register assigns requirements to decision roles. The following
table maps those roles to the staffed positions in the 12-week plan:

\begin{table}[H]
\centering
\small
\begin{tabularx}{\textwidth}{@{}p{0.30\textwidth}L@{}}
\toprule
Requirement owner & Staffed by \\
\midrule
Product engineer & Product engineer (weeks 1--12) \\
Evaluation lead & Evaluation lead (weeks 6--12, part-time weeks 1--5) \\
Document owner & Sponsor or product engineer \\
Domain editor & Sponsor (acceptance) or product engineer (implementation) \\
Independent reader & External recruit (week 11--12) \\
Security or policy owner & Product engineer (implementation) + sponsor (approval) \\
Sponsor or project owner & Sponsor \\
\bottomrule
\end{tabularx}
\caption{Role coverage for R1--R24 requirements. For the MVP, the product
engineer implements controls for all requirements. Acceptance observations
requiring domain expertise (R1, R6, R10, R18) or security approval (R11)
escalate to the sponsor. The independent reader is recruited externally for R8.}
\label{tab:role-coverage}
\end{table}

\subsection{Corpus rights}

Each corpus item carries a source path or URL, source hash, owner,
license or written permission, permitted use, redistribution status, retention
class, and consent date. The allowed statuses are
\texttt{cleared-internal}, \texttt{cleared-public},
\texttt{restricted-not-for-export}, \texttt{pending-clearance}, and
\texttt{excluded}. Pending or restricted material may support an internal
fixture only when its owner permits that use. It cannot enter a published
benchmark, external reader packet, or commercial training set without written
clearance.

\subsection{Definition of done}

The work packages have explicit exits:

\begin{enumerate}
\item \textbf{Brief, threat model, and schemas:} a complete brief, the contract
      files, a rights manifest, and threat fixtures pass validation. Code
      complete when all schemas validate, T1--T5 tests pass, and corpus rights
      manifest is populated.
\item \textbf{Protected LaTeX core:} source maps, canonical build, protected
      comparison, and CLI result replay are locked on the fixture set. Code
      complete when parser produces stable source maps, build runs in sandbox,
      and comparison produces deterministic finding IDs.
\item \textbf{Week-six checkpoint:} a second operator reproduces the core, or
      the scope is narrowed to the protected review utility. Acceptance
      complete when another engineer runs the deterministic path from
      documentation alone, or sponsor approves scope reduction.
\item \textbf{Bounded authoring extension:} plan, unit draft, preflight, and
      repair pass mutation and held-out calibration gates. Code complete when
      model invocations respect budget caps, emit valid JSON, and repair
      converges or stops within 3 cycles.
\item \textbf{Reader feasibility case:} one released packet, timed reader
      record, burden account, rights record, and all abstentions are preserved.
      Acceptance complete when independent reader completes task and burden
      calculation shows median($U_k - B_k$) $> 0$ or reports the deficit.
\item \textbf{Decision handoff:} a reproducible run, cost sensitivity, rights
      review, and proceed/revise/stop memorandum are delivered. Sponsor
      acceptance when all deliverables are received and decision is recorded.
\end{enumerate}

A work package is not complete if: (a) its mapped requirements from R1--R24
lack passing tests; (b) fixtures for this WP fail or produce unexpected
results; (c) another engineer cannot run the deliverable from documentation
alone; or (d) sponsor review identifies scope mismatch. Breaking changes are
recorded in a change log; TODOs are resolved or promoted to tracked issues.
or that the product improves expert writing. Those claims remain the province
of the named reader and the later comparative study.
